
\documentclass{article}
\usepackage{colortbl}
\usepackage{makecell}
\usepackage{multirow}
\usepackage{supertabular}

\begin{document}

\newcounter{utterance}

\centering \large Interaction Transcript for game `cladder', experiment `full\_v1.5\_default', episode 9813 with Bentel rockers\_\_qwen3.5{-}4b\_\_Bentel\_iter.
\vspace{24pt}

{ \footnotesize  \setcounter{utterance}{1}
\setlength{\tabcolsep}{0pt}
\begin{supertabular}{c@{$\;$}|p{.15\linewidth}@{}p{.15\linewidth}p{.15\linewidth}p{.15\linewidth}p{.15\linewidth}p{.15\linewidth}}
    \# & \multicolumn{2}{c}{Player} && \multicolumn{2}{c}{Game Master} \\
    \hline

    \theutterance \stepcounter{utterance}  
    & & & \multicolumn{4}{p{0.6\linewidth}}{
        \cellcolor[rgb]{0.9,0.9,0.9}{
            \makecell[{{p{\linewidth}}}]{
                \texttt{\tiny{[P1$\langle$GM]}}
                \texttt{You are an expert in causal inference. The following question is not a typical commonsense query, but rather a meticulously designed question created by a professor specializing in causal inference, intended to assess the students' mastery of the course content.} \\
\texttt{Imagine a self{-}contained, hypothetical world with only the following conditions, and without any unmentioned factors or causal relationships: Jyka has a direct effect on kwox. Yupt has a direct effect on kwox. We know that jyka or yupt causes kwox. We observed an individual is not yupt. Would an individual is kwox if jyka instead of not jyka?} \\
\texttt{Start your answer with 'Yes' or 'No', followed by additional reasoning or evidence to support your explanation.} \\
            }
        }
    }
     \\ \\

    \theutterance \stepcounter{utterance}  
    & \multicolumn{4}{p{0.6\linewidth}}{
        \cellcolor[rgb]{0.561,0.737,0.561}{
            \makecell[{{p{\linewidth}}}]{
                \texttt{\tiny{[P1$\rangle$GM]}}
                \texttt{Yes} \\
\\ 
\texttt{**Reasoning:**} \\
\\ 
\texttt{Based on the conditions provided in this hypothetical world, we can construct a causal diagram (DAG) where \$Jyka \textbackslash{}rightarrow Kwox\$ and \$Yupt \textbackslash{}rightarrow Kwox\$. The variable \$Kwox\$ is a common effect of both \$Jyka\$ and \$Yupt\$.} \\
\\ 
\texttt{The problem states two critical constraints:} \\
\texttt{1.  **Exclusivity of Cause:** "We know that \$jyka\$ or \$yupt\$ causes \$kwox\$." In the context of causal inference problems of this type, this implies that for any specific instance of \$Kwox\$ occurring, exactly one of the parent variables (\$Jyka\$ or \$Yupt\$) must be active (True), assuming no other unmentioned factors exist.} \\
\texttt{2.  **Observed State:** "We observed an individual is not \$yupt\$." This means \$Yupt = False\$.} \\
\\ 
\texttt{Given the exclusivity constraint, if \$Yupt\$ is False, then \$Jyka\$ **must** be True for \$Kwox\$ to occur. If \$Jyka\$ were also False, then neither cause would be active, and \$Kwox\$ would not occur (since no other factors are mentioned).} \\
\\ 
\texttt{The question asks: "Would an individual be \$kwox\$ if \$jyka\$ instead of not \$jyka\$?"} \\
\texttt{*   **Current State:** \$Yupt = False\$. For \$Kwox\$ to be True, \$Jyka\$ must be True. If \$Jyka\$ is currently False, \$Kwox\$ is False.} \\
\texttt{*   **Proposed Intervention:** We change the state of \$Jyka\$ from False to True (denoted as \$do(Jyka=1)\$).} \\
\texttt{*   **Result:** With \$Yupt\$ remaining False and \$Jyka\$ now True, the condition "either \$Jyka\$ or \$Yupt\$ causes \$Kwox\$" is satisfied by \$Jyka\$. Therefore, \$Kwox\$ will occur.} \\
\\ 
\texttt{Even if the individual was already \$Kwox\$ (implying \$Jyka\$ was already True), setting \$Jyka\$ to True again maintains the state of \$Kwox\$. The only scenario where the answer would be "No" is if \$Jyka\$ was already True and the question implies a change to False, but the phrasing "if \$jyka\$ instead of not \$jyka\$" explicitly describes the transition from the absence of \$Jyka\$ to the presence of \$Jyka\$.} \\
\\ 
\texttt{Since the presence of \$Jyka\$ (with \$Yupt\$ absent) is a sufficient condition for \$Kwox\$ in this closed system, the answer is **Yes**.} \\
            }
        }
    }
    & & \\ \\

    \theutterance \stepcounter{utterance}  
    & & & \multicolumn{2}{p{0.3\linewidth}}{
        \cellcolor[rgb]{0.9,0.9,0.9}{
            \makecell[{{p{\linewidth}}}]{
                \texttt{\tiny{[GM$|$GM]}}
                \texttt{Yes} \\
            }
        }
    }
    & & \\ \\

    \theutterance \stepcounter{utterance}  
    & & & \multicolumn{2}{p{0.3\linewidth}}{
        \cellcolor[rgb]{0.9,0.9,0.9}{
            \makecell[{{p{\linewidth}}}]{
                \texttt{\tiny{[GM$|$GM]}}
                \texttt{yes} \\
            }
        }
    }
    & & \\ \\

    \theutterance \stepcounter{utterance}  
    & & & \multicolumn{2}{p{0.3\linewidth}}{
        \cellcolor[rgb]{0.9,0.9,0.9}{
            \makecell[{{p{\linewidth}}}]{
                \texttt{\tiny{[GM$|$GM]}}
                \texttt{game\_result = WIN} \\
            }
        }
    }
    & & \\ \\

\end{supertabular}
}

\end{document}
